\section{问题四模型建立和求解}

\subsection{三机单弹协同模型}

设参与任务的无人机集合为 $\mathcal U_4=\{FY1,FY2,FY3\}$。对任意 $i\in\mathcal U_4$，以 $\bm x_i=(\theta_i,v_i,t_{r,i},\tau_i)^{\mathsf T}$ 描述其单弹策略。由于三架无人机互不共享航迹约束，各自的投放点和起爆点均由本机初始位置 $\bm u_{i0}$ 反算：
\begin{align}
 \bm r_{r,i}&=\bm u_{i0}+v_i\bm e(\theta_i)t_{r,i},\\
 \bm r_{e,i}&=\bm u_{i0}+v_i\bm e(\theta_i)(t_{r,i}+\tau_i)
 +(0,0,-\tfrac12g\tau_i^2)^{\mathsf T}.
\end{align}
令 $\Omega_i(\bm x_i)$ 为第 $i$ 架无人机烟幕弹对 M1 的完整圆柱遮蔽集合，则三机协同收益为 $\mu(\Omega_1\cup\Omega_2\cup\Omega_3)$。该定义允许遮蔽不连续，同时自动扣除不同无人机之间的重叠时间。

\subsection{优化模型汇总}

将三组单弹变量、烟幕寿命和几何约束统一后，问题四模型为
\begin{equation}
\begin{aligned}
&\max_{\bm x_1,\bm x_2,\bm x_3}\quad
T_4=\mu\!\left(\bigcup_{i\in\mathcal U_4}\Omega_i(\bm x_i)\right),\\[-1mm]
&\text{\rm s.t.}\quad
\left\{\begin{aligned}
&\bm x_i=(\theta_i,v_i,t_{r,i},\tau_i)^{\mathsf T},\\
&0\le\theta_i<2\pi,\quad 70\le v_i\le140,\quad t_{r,i},\tau_i\ge0,\\
&t_{e,i}=t_{r,i}+\tau_i\le t_{M1},\\
&\bm r_{r,i}=\bm u_{i0}+v_i\bm e(\theta_i)t_{r,i},\\
&\bm r_{e,i}=\bm u_{i0}+v_i\bm e(\theta_i)t_{e,i}
 +(0,0,-\tfrac12g\tau_i^2)^{\mathsf T},\\
&z_{e,i}=u_{i0,z}-\tfrac12g\tau_i^2\ge0,\\
&\bm c_i(t)=\bm r_{e,i}-(0,0,3(t-t_{e,i}))^{\mathsf T},\\
&t_{e,i}\le t\le\min\{t_{e,i}+20,t_{M1}\},\\
&D_i(t)=\displaystyle\max_{\bm p\in\mathcal T}d_i(\bm p,t),\\
&\Omega_i=\{t:D_i(t)\le10\},\qquad i\in\mathcal U_4.
\end{aligned}\right.
\end{aligned}
\label{eq:q4-complete-model}
\end{equation}
式\cref{eq:q4-complete-model} 的十二个连续变量虽可直接联合搜索，但三个单弹最优窗口在时间轴上相距较远。本文先复用问题二的单弹搜索器分别求解，再对窗口重叠部分做联合调整；当三个窗口不相交时，目标可以分解为 $T_4=\sum_i\mu(\Omega_i)$，从而显著降低搜索代价。

\subsection{求解流程与投放策略}

FY1 采用问题二发现的外向高速盆地负责早期窗口；FY2 从侧上方切入，负责中期窗口；FY3 因初始高度较低、距离真目标较近，承担晚期遮蔽。算法流程如\cref{fig:q4-flow}所示，最终策略如\cref{tab:q4-strategy}所示；三机使用不同航向和速度获得互补时间窗，而不是重复覆盖同一阶段。

\begin{figure}[!htbp]
 \centering
 \includegraphics[width=0.58\textwidth]{fig_flow_q4.pdf}
 \caption{问题四三机单弹协同优化流程}
 \label{fig:q4-flow}
\end{figure}

\begin{table}[!htbp]
 \centering
 \caption{问题四三架无人机协同投放策略}
 \label{tab:q4-strategy}
 \resizebox{\textwidth}{!}{
 \begin{tabular}{ccccccc}
  \toprule[1.5pt]
  无人机 & 航向角/$^{\circ}$ & 速度/(m/s) & 投放时刻/s & 引信延时/s & 起爆点/m & 单弹时长/s \\
  \midrule[1pt]
  FY1 & $5.112023$ & $140.000000$ & $0.878620$ & $0.053112$ & $(17929.924,11.623,1799.986)$ & $4.588062$ \\
  FY2 & $306.190938$ & $136.853230$ & $8.551220$ & $3.985720$ & $(13013.095,15.321,1322.159)$ & $3.836164$ \\
  FY3 & $122.474312$ & $92.960170$ & $31.732660$ & $7.645810$ & $(4034.530,88.225,413.554)$ & $2.232822$ \\
  \bottomrule[1.5pt]
 \end{tabular}}
\end{table}

三枚弹的有效区间依次为
\begin{equation}
 [0.931732,5.519794],\quad[17.730838,21.567002],\quad[41.457952,43.690774].
\end{equation}
三个窗口两两不交，因此联合时长等于三者之和：
\begin{equation}
 T_4=4.588062160+3.836163637+2.232821992
 =\boxed{10.657047788\ \mathrm{s}}.
\end{equation}
早、中、晚三段分工说明多机协同的主要收益来自空间可达性带来的时间窗互补，而非在同一时刻堆叠更多烟幕。具体投放点与起爆点已写入 \texttt{result2.xlsx}。

三架无人机的单弹贡献与联合时间窗如\cref{fig:q4-diagnostics}所示。FY1、FY2 和 FY3 分别形成 $4.588$、$3.836$ 和 $2.233\,\mathrm{s}$ 的早、中、晚窗口，三段在时间轴上互不重叠，因此黑色联合窗口长度等于三段时长之和 $10.657\,\mathrm{s}$，与式中终算结果一致。

\begin{figure}[!htbp]
 \centering
 \includegraphics[width=0.96\textwidth]{q4_coordination_diagnostics.pdf}
 \caption{问题四单弹贡献与三机错峰遮蔽窗口}
 \label{fig:q4-diagnostics}
\end{figure}

\subsection{结果复核}

终算采用每个圆周 $5760$ 个角向点，时间扫描步长为 $0.005\,\mathrm{s}$，并用事件根修正所有进出时刻。三枚烟幕弹均在导弹到达假目标前起爆，起爆高度均为正；FY1、FY2、FY3 的有效窗口之间至少相隔 $12.21\,\mathrm{s}$，故轻微执行误差不会引起跨机窗口重叠。由粗算的 $10.657051568\,\mathrm{s}$ 加密至终算后仅变化 $3.78\times10^{-6}\,\mathrm{s}$，结果具有良好的离散收敛性。

\endinput
\section{问题四模型建立和求解}

参考文献对一篇正式论文来说是必不可少的，在建模中重要的参考文献当然应该列出。\LaTeX{}在这方面的功能也是十分强大的，下面介绍一个比较简单的参考文献制作方法。有兴趣的可学习 \verb|bibtex| 或 \verb|biblatex| 的使用。

\LaTeX{}的入门书籍可以看《\LaTeX{}入门》\cite{liuhaiyang2013latex}。这是一个简单的引用，用 \verb|\cite{bibkey}| 来完成。

本模板还提供了 \verb|\upcite{bibkey}| 和 \verb|\supercite{bibkey}| 命令用来引用，渲染效果是上标。
